%% GENERATED by python/paper/build.py from results/ -- do not edit by hand.
%% Rebuild:  .venv\Scripts\python.exe python\paper\build.py --only US_CRRA_PensChars_vectorX
%% Source:   results/shocks/US_shocks.csv
\begin{table}[!htb]
\centering
\begin{threeparttable}
\caption{Does CRRA matter for the effect of pension design ($\theta$) in US -- 2020 (vector $X_i$ calibration)}
\label{table:US:CRRA:pensChars_vectorX}
\renewcommand{\arraystretch}{1.25}
\begin{tabularx}{.9\textwidth}{YYYYY}
\toprule
\textbf{Scenario} & \textbf{CRRA} ($\rho$) & \textbf{Tax rate} & \textbf{Savings rate} & \textbf{Avg. workweek} \\
\midrule 
 & 0.5 & 14.43\% & 15.91\% & 39.39 \\
Baseline & 1.0 & 14.43\% & 15.37\% & 39.39 \\
 & 2.0 & 14.43\% & 14.85\% & 39.39 \\[.5em]\hline\\[-.75em]
 & 0.5 & 20.23\% & $-2.18$ p.p. & 36.80 \\
$\theta = 0$ & 1.0 & 18.52\% & $-0.93$ p.p. & 37.52 \\
 & 2.0 & 15.21\% & $+0.10$ p.p. & 38.38 \\[.5em]\hline\\[-.75em]
 & 0.5 & 11.21\% & $+1.54$ p.p. & 40.55 \\
$\theta = 1$ & 1.0 & 12.28\% & $+0.65$ p.p. & 40.13 \\
 & 2.0 & 13.72\% & $+0.06$ p.p. & 39.79 \\[.5em]\hline\\[-.75em]
\bottomrule
\end{tabularx}
\begin{tablenotes}
\footnotesize
\item Every $\rho$ is separately calibrated: the tax rate is a target and the workweek is normalised against each $\rho$\textquotesingle s own baseline, so both are common to the baseline rows, while the baseline savings rate is a prediction ($\beta$ is identified by the interest rate) and falls with $\rho$. The savings rate is savings relative to GDP; the baseline rows report its level and every scenario row the change against the baseline at the same $\rho$, in percentage points. Vector-$X_i$ calibration: $X_i$ is identified from relative hours, which are data here, and the level of $ar h$ is then not identified --- only its ratio to the baseline is. $eta$, $\omega$, the tax rate and the savings rate are common to the two variants; what differs is anything defined through $\eta$ or $X$.
\end{tablenotes}
\end{threeparttable}
\end{table}
